\section{Summary and outlook}
\label{sec:summary}
We have developed a systematic approach for matching from the gradient-flow scheme to the $\overline{\rm MS}$ scheme for Wilson-line operators in the small flow-time limit. This approach extends to various applications, including quasi-PDFs, gluonic correlators that appear in quarkonium decay and production in the framework of pNRQCD and spin-dependent potentials in terms of chromomagnetic field insertions into a Wilson loop.

Within the one-dimensional auxiliary-field formalism, the matching from the gradient-flow scheme to the $\overline{\rm MS}$ scheme for the Wilson-line operators discussed in this work can be reduced to the matching for the corresponding local current operators in the small flow-time limit.
We match these local current operators from the gradient-flow scheme to the $\overline{\rm MS}$ scheme by comparing perturbative computations in these two schemes.
Great simplification for the matching procedure can be achieved by using small flow-time expansion and setting external momenta to be zero whenever possible.
In this way, the well-developed method of matching calculation in HQET can be applied. After completing the matching for the local current operators, one automatically obtains the matching coefficients for the corresponding Wilson-line operators in the small flow-time limit,
because the combinations of the renormalized local current operators at different spacetime do not introduce further UV divergences, apart from the subtracted linear divergences, due to the multiplicative renormalizability of the Wilson-line operators.


The matching for the local current operators can be further decomposed as the matching for the fields and vertices due to the multiplicative renormalizability of the local current operators.
We thus compute the matching coefficients at one-loop level for auxiliary ``heavy” auxiliary field $h_v$,  massless quark field $\psi$, $\psi h_v$ vertex,  gluon field $gA$ and the vertices of the local current operators  $gF^{\mu\nu}_{\|\perp}h_v,\, gF^{\mu\nu}_{\perp\perp}h_v$ and $g\bar{h}_vF^{\mu\nu}_{\perp\perp}h_v$. The one-loop result of the linearly divergent ``mass” correction $\delta m$ is also obtained. Provided with these results, we give the matching coefficients for the local current operators and the Wilson-line operators in the small flow-time limit. 

It is worth mentioning that this method is also applicable for matching from lattice regularization schemes to the $\overline{\rm MS}$ scheme in the small lattice spacing limit, albeit with added complexity due to UV and IR divergence treatment \cite{Constantinou:2017sej}.



In the appendix, we have shown that, at one-loop level, for the case of hadronic matrix element defined in eq.~(\ref{eq:quarkhmte}) at zero external momenta, the finite flow time effect is very small as long as the flow radius is smaller than physical distance $z$, which is usually satisfied in lattice gradient flow computations. 


With the systematic approach we have developed for the matching from the gradient-flow scheme to the $\overline{\rm MS}$ scheme for the Wilson-line operators, 
 it is straightforward to extend the current study to two-loop order, which would be more helpful for lattice gradient flow computations of matrix elements constructed using the off-lightcone Wilson-line operators. We leave it for future research.
 
\subsubsection*{Note added}
Before submission of this work, we noticed that the one-loop result of the matching coefficient for the gluonic correlator $G_B$ was recently published in ref.~\cite{Altenkort:2023eav},  which was the same as our one-loop result of  $c_F^2(t, \mu)$, because the matching calculation for the gluonic correlator $G_B$ could be reduced to the matching calculation for the local current operator $g\bar{h}_vF^{\mu\nu}_{\perp\perp}h_v$.

